% Part: turing-machines
% Chapter: machines-computations
% Section: representing-tms

\documentclass[../../../include/open-logic-section]{subfiles}

\begin{document}

\olfileid[pt]{tur}{mac}{rep}
\olsection{Representando Máquinas de Turing}

\begin{explain}
As máquinas de Turing podem ser representadas visualmente por \emph{diagramas de estado}. 
Os diagramas são compostos por
células de estado conectadas por setas. Sem surpresa, cada célula de estado representa
um estado da máquina. Cada seta representa uma instrução que pode ser
executada a partir desse estado, com as especificações da instrução escritas acima
ou abaixo da seta apropriada. Considere a seguinte máquina, que tem 
apenas dois estados internos, $q_0$ e $q_1$, e uma instrução:
\[
\begin{tikzpicture}[->,>=stealth',shorten >=1pt,auto,node distance=2.8cm,
                    semithick]
  \tikzstyle{every state}=[fill=none,draw=black,text=black]

  \node[initial,state]   (A)              {$q_0$};
  \node[state]   (B) [right of=A] {$q_1$};

  \path (A) edge  node {\TMtrans{\TMblank}{\TMstroke}{\TMright}} (B);
\end{tikzpicture}
\]
Lembre-se de que a máquina de Turing tem uma cabeça de leitura/escrita e uma fita com
a entrada escrita nela. A instrução pode ser lida como \emph{se
lendo um~$\TMblank$ no estado $q_0$, escreva um~$\TMstroke$, mova-se à direita,
e vá para o estado $q_1$}. Isto é equivalente à função de
transição que mapeia $\tuple{q_0, \TMblank}$ em $\tuple{q_1, \TMstroke,
\TMright}$.
\end{explain}

\begin{ex}
\emph{Máquina Par}: A seguinte máquina de Turing para se, e somente
se, houver um número par de $\TMstroke$ na fita (sob a
suposição de que todos os $\TMstroke$ vêm antes do primeiro $\TMblank$ na
fita).
\[
\begin{tikzpicture}[->,>=stealth',shorten >=1pt,auto,node distance=2.8cm,
                    semithick]
  \tikzstyle{every state}=[fill=none,draw=black,text=black]

  \node[initial,state]         (A)              {$q_0$};
  \node[state]         (B) [right of=A] {$q_1$};

  \path (A) edge [bend left] node {\TMtrans{\TMstroke}{\TMstroke}{\TMright}} (B)
        (B) edge [loop above] node {\TMtrans{\TMblank}{\TMblank}{\TMright}} (B)
            edge [bend left] node {\TMtrans{\TMstroke}{\TMstroke}{\TMright}} (A);
\end{tikzpicture}
\]

O diagrama de estado corresponde à seguinte função de transição:
\begin{align*}
\delta(q_0, \TMstroke) & = \tuple{q_1, \TMstroke, \TMright},\\
\delta(q_1, \TMstroke) & = \tuple{q_0, \TMstroke, \TMright},\\
\delta(q_1, \TMblank)  & = \tuple{q_1, \TMblank, \TMright}
\end{align*}
\end{ex}

\begin{explain}
A máquina acima para apenas quando a entrada é um número par de traços.
Caso contrário, a máquina (teoricamente) continua a operar indefinidamente. 
Para qualquer máquina e entrada, é possível rastrear as 
\emph{configurações} da máquina a fim de determinar a saída. 
Daremos uma definição formal de configurações mais adiante. Por enquanto,
podemos pensar intuitivamente nas configurações como uma série de diagramas que mostram o 
estado da máquina em qualquer instante durante a operação.
As configurações mostram o conteúdo da fita, o estado da máquina e a
posição da cabeça de leitura/escrita.

Vamos rastrear as configurações da máquina par se ela for
iniciada com uma entrada de quatro $\TMstroke$. Nesse caso, esperamos que a
máquina pare.  Em seguida, executaremos a máquina com uma entrada de três
$\TMstroke$, na qual a máquina rodará para sempre.

A máquina inicia no estado~$q_0$, examinando o~$\TMstroke$ mais à esquerda.
Podemos representar o estado inicial da máquina da seguinte forma:
\[
\TMendtape \TMstroke_0 \TMstroke \TMstroke \TMstroke \TMblank \ldots
\]
A configuração acima é direta. Como se pode ver, a
máquina inicia no estado um, examinando o~$\TMstroke$ mais à esquerda. Isso
é representado por um subscrito do nome do estado no
primeiro~$\TMstroke$. A instrução aplicável neste ponto é
$\delta(q_0, \TMstroke) = \tuple{q_1, \TMstroke, \TMright}$, e assim
a máquina move-se à direita na fita e muda para o estado~$q_1$.
\[
\TMendtape \TMstroke \TMstroke_1 \TMstroke \TMstroke \TMblank \ldots
\]
Como a máquina está agora no estado~$q_1$ examinando um~$\TMstroke$, temos que
``seguir'' a instrução $\delta(q_1, \TMstroke) = \tuple{q_0,
  \TMstroke, \TMright}$. Isso resulta na configuração
\[
\TMendtape \TMstroke \TMstroke \TMstroke_0 \TMstroke \TMblank \ldots
\]
À medida que a máquina continua, as regras são aplicadas novamente na mesma
ordem, resultando nas duas configurações a seguir:
\[
\TMendtape \TMstroke \TMstroke \TMstroke \TMstroke_1 \TMblank \ldots
\]
\[
\TMendtape \TMstroke \TMstroke \TMstroke \TMstroke \TMblank_0 \ldots
\]
A máquina está agora no estado~$q_0$ examinando um~$\TMblank$. Com base no
diagrama de transição, podemos ver facilmente que não há nenhuma instrução a ser
executada, e portanto a máquina parou. Isso significa que a entrada
foi aceita.

Suponha que, em seguida, iniciemos a máquina com uma entrada de três
$\TMstroke$. As primeiras configurações são semelhantes, pois as mesmas
instruções são executadas, com apenas uma pequena diferença na entrada
da fita:
\[
\TMendtape \TMstroke_0 \TMstroke \TMstroke \TMblank \ldots
\]
\[
\TMendtape \TMstroke \TMstroke_1 \TMstroke \TMblank \ldots
\]
\[
\TMendtape \TMstroke \TMstroke \TMstroke_0 \TMblank \ldots
\]
\[
\TMendtape \TMstroke \TMstroke \TMstroke \TMblank_1 \ldots
\]
A máquina já percorreu todos os $\TMstroke$ e está
lendo um~$\TMblank$ no estado~$q_1$. Como mostrado no diagrama, há
uma instrução da forma $\delta(q_1, \TMblank) =\tuple{q_1,
\TMblank, \TMright}$. Como a fita está preenchida com $\TMblank$
indefinidamente para a direita, a máquina continuará a executar esta
instrução \emph{para sempre}, permanecendo no estado~$q_1$ e movendo-se cada vez
mais à direita. A máquina nunca para, e não aceita
a entrada.
\end{explain}

\begin{explain}
É importante observar que nem todas as máquinas param. Se parar
significa que a máquina fica sem instruções a executar, então
podemos criar uma máquina que nunca para, simplesmente garantindo que haja
uma seta de saída para cada símbolo em cada estado. A máquina par pode
ser modificada para rodar indefinidamente adicionando uma instrução para examinar um
$\TMblank$ em~$q_0$.
\end{explain}

\begin{ex}
\[
\begin{tikzpicture}[->,>=stealth',shorten >=1pt,auto,node distance=2.8cm,
                    semithick]
  \tikzstyle{every state}=[fill=none,draw=black,text=black]

  \node[initial,state] (A)              {$q_0$};
  \node[state]         (B) [right of=A] {$q_1$};

  \path
  (A) edge [bend left] node {\TMtrans{\TMstroke}{\TMstroke}{\TMright}} (B)
      edge [loop above] node {\TMtrans{\TMblank}{\TMblank}{\TMright}} (A)
  (B) edge [loop above] node {\TMtrans{\TMblank}{\TMblank}{\TMright}} (B)
      edge [bend left] node {\TMtrans{\TMstroke}{\TMstroke}{\TMright}} (A);
\end{tikzpicture}
\]
\end{ex}

\begin{explain}
As tabelas de máquina são outra forma de representar máquinas de
Turing. As tabelas de máquina têm o alfabeto da fita exibido no
eixo~$x$, e o conjunto de estados da máquina ao longo do eixo~$y$. Dentro da
tabela, na interseção de cada estado e símbolo, está escrito o
resto da instrução---o novo estado, o novo símbolo e a direção do
movimento. As tabelas de máquina facilitam determinar em que estado, e
para qual símbolo, a máquina para. Sempre que há uma lacuna na
tabela, há um ponto possível para a máquina parar. Ao contrário dos diagramas
de estado e dos conjuntos de instruções, em que os pontos nos quais a máquina
para nem sempre são imediatamente óbvios, quaisquer pontos de parada são
rapidamente identificados ao encontrar as lacunas na tabela de máquina.
\end{explain}

\begin{ex}
A tabela de máquina da máquina par é:
\[
\centering
\begin{tabular}{lllll}
\cline{2-4}
\multicolumn{1}{l|}{}      & \multicolumn{1}{c|}{$\TMblank$}                
& \multicolumn{1}{c|}{$\TMstroke$}     & \multicolumn{1}{c|}{$\TMendtape$}   &  \\ \cline{1-4}
\multicolumn{1}{|l|}{$q_0$} & \multicolumn{1}{l|}{}                          
& \multicolumn{1}{l|}{\TMtrans{\TMstroke}{q_1}{\TMright}} & 
\multicolumn{1}{l|}{\phantom{\TMtrans{\TMstroke}{q_1}{\TMright}}} &  \\ \cline{1-4}
\multicolumn{1}{|l|}{$q_1$} & \multicolumn{1}{l|}{\TMtrans{\TMblank}{q_1}{\TMright}} 
& \multicolumn{1}{l|}{\TMtrans{\TMstroke}{q_0}{\TMright}} & 
\multicolumn{1}{l|}{\phantom{\TMtrans{\TMstroke}{q_1}{\TMright}}} &  \\ \cline{1-4}
\end{tabular}
\]
Como podemos ver, a máquina para ao examinar um~$\TMblank$ no estado~$q_0$.
\end{ex}

\begin{explain}
Até aqui consideramos apenas máquinas que leem e aceitam entradas.
No entanto, as máquinas de Turing têm a capacidade de tanto ler quanto escrever. Um
exemplo de tal máquina (embora haja muitos, muitos exemplos) é
um \emph{duplicador}. Um duplicador, quando iniciado com um bloco de~$n$
$\TMstroke$ na fita, produz como saída um bloco de~$2n$ $\TMstroke$.
\end{explain}

\begin{ex}
  \ollabel{ex:doubler} Antes de construir uma máquina duplicadora, é
  importante elaborar uma \emph{estratégia} para resolver o problema.
  Como a máquina (tal como a formulamos) não consegue lembrar quantos
  $\TMstroke$ já leu, precisamos encontrar uma maneira de
  rastrear todos os $\TMstroke$ na fita. Uma dessas maneiras é
  separar a saída da entrada com um~$\TMblank$. A máquina
  pode então apagar o primeiro $\TMstroke$ da entrada, percorrer
  o resto da entrada, deixar um $\TMblank$, e escrever dois novos
  $\TMstroke$.  A máquina então voltará e encontrará o segundo
  $\TMstroke$ na entrada, e duplicará esse também. Para cada
  $\TMstroke$ de entrada, ela escreverá dois $\TMstroke$ de saída.  Ao
  apagar a entrada conforme a máquina avança, podemos garantir que nenhum
  $\TMstroke$ seja perdido ou duplicado duas vezes. Quando toda a entrada é
  apagada, restarão $2n$ $\TMstroke$ na fita. O diagrama de
  estado da máquina de Turing resultante está representado na
  \olref{fig:doubler}.
  \begin{figure}
\[
\begin{tikzpicture}[->,>=stealth',shorten >=1pt,auto,node distance=2.8cm,
                    semithick]
  \tikzstyle{every state}=[fill=none,draw=black,text=black]

  \node[initial,state] (1)              {$q_0$};
  \node[state]         (2) [right of=1] {$q_1$};
  \node[state]         (3) [right of=2] {$q_2$};
  \node[state]         (4) [below of=3] {$q_3$};
  \node[state]         (5) [left of=4]  {$q_4$};
  \node[state]         (6) [left of=5]  {$q_5$};

  \path (1) edge node {\TMtrans{\TMstroke}{\TMblank}{\TMright}} (2)
    (2) edge [loop above] node {\TMtrans{\TMstroke}{\TMstroke}{\TMright}} (2)
      edge node {\TMtrans{\TMblank}{\TMblank}{\TMright}} (3)
    (3) edge [loop above] node {\TMtrans{\TMstroke}{\TMstroke}{\TMright}} (3)
        edge node {\TMtrans{\TMblank}{\TMstroke}{\TMright}} (4)
    (4) edge [loop below] node {\TMtrans{\TMblank}{\TMstroke}{\TMleft}} (4)
        edge node {\TMtrans{\TMstroke}{\TMstroke}{\TMleft}} (5)
    (5) edge [loop below]  node {\TMtrans{\TMstroke}{\TMstroke}{\TMleft}} (5)
        edge              node {\TMtrans{\TMblank}{\TMblank}{\TMleft}} (6)
    (6) edge [loop below] node {\TMtrans{\TMstroke}{\TMstroke}{\TMleft}} (6)
        edge              node {\TMtrans{\TMblank}{\TMblank}{\TMright}} (1);
\end{tikzpicture}
\]
\caption{Uma máquina duplicadora}
\ollabel{fig:doubler}
\end{figure}
\end{ex}

\begin{prob}
Escolha uma entrada arbitrária e rastreie as configurações da
máquina duplicadora em \olref[tur][mac][rep]{ex:doubler}.
\end{prob}

\begin{prob}
Projete uma máquina de Turing com alfabeto $\{\TMendtape,\TMblank, A, B\}$
que aceite, isto é, pare, qualquer cadeia de $A$ e $B$ na qual o
número de $A$ seja igual ao número de $B$ \emph{e} todos os
$A$ precedam todos os $B$, e que rejeite, isto é, não pare, qualquer
cadeia em que o número de $A$ não seja igual ao número de~$B$, ou
os $A$ não precedam todos os~$B$. (Por exemplo, a máquina deve aceitar
$AABB$, e $AAABBB$, mas rejeitar tanto $AAB$ quanto $AABBAABB$.)
\end{prob}

\begin{prob}
Projete uma máquina de Turing com alfabeto $\{\TMendtape,\TMblank, A, B\}$
que receba como entrada qualquer cadeia $\alpha$ de $A$ e $B$ e
a duplique para produzir uma saída da forma $\alpha\alpha$. (Por exemplo,
a entrada $ABBA$ deve resultar na saída $ABBAABBA$).
\end{prob}

\begin{prob}
\emph{Alfabética?:} Projete uma máquina de Turing com alfabeto
$\{\TMendtape,\TMblank, A, B\}$ que, ao receber como entrada uma sequência
finita de $A$ e $B$, verifique se todos os $A$ aparecem à
esquerda de todos os $B$ ou não. A máquina deve deixar a cadeia de
entrada na fita, e ou parar, se a cadeia for
``alfabética'', ou rodar para sempre, se a cadeia não for.
\end{prob}

\begin{prob}
\emph{Ordenador alfabético:} Projete uma máquina de Turing com alfabeto
$\{\TMendtape,\TMblank, A, B\}$ que receba como entrada uma sequência finita
de $A$ e $B$ e os reorganize de modo que todos os $A$ fiquem à
esquerda de todos os~$B$. (Por exemplo, a sequência $BABAA$ deve tornar-se a
sequência $AAABB$, e a sequência $ABBABB$ deve tornar-se a sequência
$AABBBB$).
\end{prob}

\end{document}
